\section{Measurement Framework and Problem Formulation}
\label{sec:framework}

\begin{figure}[tb]
\centering
\includegraphics[width=\columnwidth]{framework.pdf}
\caption{Measurement framework, read top to bottom. The optical path supplies
the reference, the inertial path the window-demeaned specific force, and the
twin the synthetic streams. Neither path observes the initial velocity inside
a window. In the twin block, $\boldsymbol\Sigma$, ${\bf b}\!\sim\!{\rm RW}$,
$Q$, $\delta t$ and ${\bf g}_B$ are the noise covariance, bias random walk,
quantization, sampling jitter and board-frame gravity. The red dashed line
marks that the held-out recording does not flow back into fitting or
cross-validation.}
\label{fig:framework}
\end{figure}

\subsection{Measurand and Observable Information}
Let $B$, $C$, and $I$ denote the board, left-camera, and USB-IMU frames. A
rotation ${\bf R}_{XY}$ maps the coordinates of a vector from frame $Y$ to
frame $X$; in particular, ${\bf R}_{CI}$ maps IMU vectors into the camera
frame. The
measurand is the time-varying camera pose
\begin{equation}
 {\bf T}_{BC}(t)=
 \begin{bmatrix}{\bf R}_{BC}(t)&{\bf p}_{BC}(t)\\{\bf 0}^{\top}&1\end{bmatrix}
 \in SE(3),
\end{equation}
where ${\bf p}_{BC}$ is the camera-center position expressed in $B$. For two
times $t_a<t_b$ with $\Delta t_{ab}=t_b-t_a$ between $0.6h$ and $1.5h$ around
the nominal horizon $h=3$~s (the sweep that selects $h$ is reported in
Section~\ref{sec:limit}), the learning target is the translation increment
expressed in the start camera frame,
\begin{equation}
 \Delta{\bf p}_{ab}={\bf R}_{BC}^{\top}(t_a)
 [ {\bf p}_{BC}(t_b)-{\bf p}_{BC}(t_a) ] . \label{eq:target}
\end{equation}
This definition removes the arbitrary board orientation while retaining
metric scale. Although the reference is a full $SE(3)$ pose, the learned
output in this study is the three-dimensional translation increment and its
integrated translation trajectory; the network does not claim a learned
six-degree-of-freedom pose. Reference orientation is used only in the
translation-isolated trajectory analysis, and the pure-inertial deployment
input remains the IMU stream alone.

The IMU reports specific force $\tilde{\bf f}_I$ and angular rate
$\tilde{\boldsymbol\omega}_I$ at approximately 200~Hz. Camera frames are
timestamped at 10--30~Hz. The constant rotation ${\bf R}_{CI}$ is estimated
in a dedicated rigidity recording, not from validation or test trajectories.

\label{sec:observable}
Writing the measurand as a double integral of acceleration separates it into
two parts with different information status,
\begin{equation}
 \Delta{\bf p}_{ab}=
 \underbrace{{\bf v}(t_a)\,\Delta t_{ab}}_{\text{initial condition}}
 +\underbrace{\int_{t_a}^{t_b}\!\!\int_{t_a}^{\tau'}{\bf a}(\tau)\,
 {\rm d}\tau\,{\rm d}\tau'}_{\text{inertially measured}},
 \label{eq:decomp}
\end{equation}
where ${\bf v}(t_a)$ and ${\bf a}(\tau)$ are respectively the camera-center
velocity at $t_a$ and its acceleration, both resolved in the fixed start frame
$C_a$. The IMU measures the specific force at its own sensing point, so
${\bf a}$ follows from ${\bf R}_{CI}\tilde{\bf f}_I$ after rotation into
$C_a$, addition of gravity, and the lever-arm term of
Section~\ref{sec:learning}; the second term is therefore determined by the
inertial measurement up to bias, scale, and the gravity that an attitude error
leaks into the integrand. The first term is an \emph{initial condition}. Adding
a constant velocity ${\bf v}_0$ to a trajectory,
${\bf p}(t)\mapsto{\bf p}(t)+{\bf v}_0t$, leaves $\ddot{\bf p}$ and ${\bf R}$,
and hence every specific-force and angular-rate sample, unchanged. Therefore
${\bf v}(t_a)$ is not a function of the inertial record on the context window
$[t_a-c,\,t_b+c]$, where $c>0$ is the context duration on each side of the
target interval, unless a velocity boundary condition is available inside it.
A windowed estimator can only replace ${\bf v}(t_a)$ by a prior learned from
the training motion; under squared-error loss the optimal data-independent
prior is the training mean, which is close to zero for a handheld instrument
that reverses direction. Its predictions should then be shrunk toward zero,
and Section~\ref{sec:limit} tests this prediction.

Two acquisition properties can reduce this ambiguity. A longer context does
not make ${\bf v}(t_a)$ observable; it only supplies more motion history from
which a learned prior may infer low-frequency velocity. A correctly
identified stationary epoch $t_0$ instead supplies the boundary condition
${\bf v}(t_0)={\bf0}$, so that
$\hat{\bf v}(t_a)=-\int_{t_a}^{t_0}\hat{\bf a}\,{\rm d}\tau$ can be estimated
from the inertial stream, with $\hat{\bf a}$ computed from the estimated
attitude. That estimate remains limited by attitude, bias and
integration errors, and its usable interval is limited by the same gravity
leakage that affects the second term of (\ref{eq:decomp}). Sections
\ref{sec:limit} and \ref{sec:pause} examine both effects for this
instrument.

\subsection{Data Partitioning and Outputs}
The framework separates four information classes: \emph{calibration} (camera
intrinsics, time delay, ${\bf R}_{CI}$, board scale, IMU noise parameters);
\emph{reference} (cleaned planar-PnP pose while the board is visible and
verified stationary); \emph{learning} (complete training recordings, the
rotation-rich rigidity recordings that are never scored, and synthetic-train
sequences); and \emph{evaluation} (the validation recordings for selection in
the fixed split, leave-one-session-out folds over all non-test recordings, and
one held-out real-test recording). Test data do not parameterize the twin,
labels or gradient fitting; Section~\ref{sec:experiments} distinguishes the
initial selection protocol from subsequent exploratory diagnostics on this
recording. The boundary matters here because using the gyroscope to modify
PnP labels that are then predicted from the same gyroscope biases the reported
inertial accuracy upward, as Section~\ref{sec:imu} quantifies.

\paragraph{Outputs.}
The system produces: 1) a characterized camera-pose reference with stated
limits; 2) paired ideal/noisy synthetic sensor streams with exact pose;
3) a displacement stream anchored at the pair start, with empirical
uncertainty from a model ensemble, whose end point is the local
pseudo-measurement; 4) overlapping displacement edges---one per window, or
one per reference frame from the stream---that can be fused into a
trajectory; and 5) tests of whether the residual behaves as the unobserved
initial velocity predicts under the chosen window. The learned trajectory is
an estimator output, not a new ground-truth source.
